\documentclass[12pt]{article}

\input{0a - preambule}

\title{Pourquoi le Théorème du Produit Nul (TPN) ?}

\begin{document}

\maketitle

\textit{"si un produit est nul, alors l’un de ses facteurs est nul"} sens direct (1); \\
\textit{"si un des facteurs d’un produit est nul, alors le produit est nul"} réciproque (2).

Ces deux phrases se résument en langage mathématique ainsi:
\[
ab = 0 \;\iff a=0 \text{ ou } b=0
\]

Prouvons-le. Pour cela, il faut démontrer ces deux points suivants:
\begin{gather}   
    ab = 0 \ \Rightarrow \ a=0 \text{ ou } b=0 \\
    a=0 \text{ ou } b=0 \ \Rightarrow \ ab = 0
\end{gather}

La réciproque (2) est immédiate à démontrer car zéro est absorbant. 

Il reste donc à démontrer le sens direct (1).

Supposons que $a \neq 0$. \\
$a$ admet donc un inverse qui est $a^{-1}$, tel que $a^{-1}a = 1$. On peut donc écrire
\begin{gather*}
    ab = 0 \ \Rightarrow \ a^{-1}ab = a^{-1}\cdot0 \\
    (a^{-1}a)b = 0 \\
    1 \cdot b = 0 \\
    b = 0
\end{gather*}

On  raisonne de la même manière en supposant $b \neq 0$.

On en déduit donc
\[
ab = 0 \;\Rightarrow\; a=0 \text{ ou } b=0
\]

\end{document}
